\clearpage
\VerdierPhysicalPageAnchor{0023}
$\mathcal A$ 有足够多的内射对象。因此，全导出函子尤其能够重构
“$F$ 的右超导出函子”，但它提供的信息更多。例如，若有阿贝尔范畴之间的
加性函子 $G:\mathcal B\to\mathcal C$，就可以复合全导出函子：

\VerdierSemanticAnchor{intro-3-derived-composition}
\[
\mathrm R G\mathbin{\cdot}\mathrm R F:
\mathbf D^*(\mathcal A)\longrightarrow\mathbf D(\mathcal C).
\]
在通常情形下，函子 $\mathrm R G\mathbin{\cdot}\mathrm R F$ 与 $GF$ 的
全导出函子典范同构：

\VerdierSemanticAnchor{intro-3-derived-composition-isomorphism}
\[
\mathrm R(GF)\xrightarrow{\sim}
\mathrm R G\mathbin{\cdot}\mathrm R F,
\]
这一复合公式用我们的语言表达了\emph{复合函子的谱序列}，而且更为精确：
它不只给出 $\mathrm R GF(X)$ 的上同调对象的若干信息
（$X$ 是 $\mathbf D^*(\mathcal A)$ 的对象），而是完整描述产生这些上同调
对象的那个“在拟同构意义下”的复形。知道这个复形，例如就可以再次与另一
个全导出函子复合。

\VerdierSemanticAnchor{intro-4}
\noindent\textbf{4.}\quad
一般而言，在这一形式体系中，我们必须直接处理“在拟同构意义下”的复形，
并谨慎避免只用它们的上同调对象来取代它们。（尽管计算或论证终了时所寻求
的“不变量”通常是复形的上同调对象，但并非总是如此：谱序列的微分也是值得
关注的“不变量”。）

\VerdierSemanticAnchor{intro-4-derived-relations}
在拟同构意义下的复形层面、全导出函子层面，或更一般地在导出范畴之间的
函子层面，可以建立各种关系（同构关系、用于对偶型定理的伴随关系等）。
通过纯粹自动的过程（谱序列），这些关系会给出复形上同调对象、这些对象所
带的自然滤过等方面的信息。

\VerdierSemanticAnchor{intro-4-technical-intermediaries}
即使只关心传统导出函子的值，也就是只关心上同调对象，导出范畴中的关系
仍常常是研究这些函子并为其表述某些重要关系时不可缺少的技术中介。

\VerdierSemanticAnchor{intro-4-simpler-relations}
最后，一个根本优点是：在导出范畴中建立的关系，无论在形式上还是实质上，
都比由此推得的上同调对象之间的关系更简单。
